\clearpage
\fancyhead[L]{\small INHALT.}
\SourceLeaf{669}
\begingroup
\raggedright
\noindent\hangindent=1em \textsc{Delambre}: Tables écliptiques .des; satellites de Jupiter, d'après la théorie de M. le marquis de \textsc{Laplace}, etc. Paris 1817 \mbox{.\ .\ .\raisebox{-.3ex}{.}\ .\ .\ .\ .\ .\ .} \hfill\mbox{1818 März 19}\quad\mbox{Seite 603}\par\smallskip
\noindent\hangindent=1em \textsc{Prosperin}: Bahnen der Nebenplaneten um die Sonne etc.\ .\ .\ . Nova acta. regiae societatis scientiarum Upsaliensis \mbox{.\ .\ .\ .\ .\ .\ .\ .\ .\ .\ .\ .} \hfill\mbox{1818 Mai 2}\quad\mbox{--- 604}\par\smallskip
\noindent\hangindent=1em \textsc{Bode}: Astronomisches Jahrbuch für das Jahr 1820 \mbox{.\ \raisebox{.25ex}{-}\ .\ .\ .\ .\ .\ .\ .\ .} \hfill\mbox{1818 Mai 21}\quad\mbox{--- 605}\par\smallskip
\noindent\hangindent=1em Connaissance des tems etc.\ .\ .\ . pour l'an 1820. \mbox{.\ .\ .\ .\ .\ .\ .\ .\ .\ .\ .} \hfill\mbox{1819 März 11}\quad\mbox{--- 609}\par\smallskip
\noindent\hangindent=1em \textsc{Bode}: Astronomisches Jahrbuch für das Jahr 1821 \mbox{.\ .\ .\ .\ .\ .\ .\ .\ .\ .} \hfill\mbox{1819 Mai 20}\quad\mbox{--- 611}\par\smallskip
\noindent\hangindent=1em Connaissance des tems etc.\ .\ .\ . pour l'an 1821 \mbox{.\ .\ .\ .\ .\ .\ .\ .\ .\ .\ .} \hfill\mbox{1819 Nov. 29}\quad\mbox{--- 615}\par\smallskip
\noindent\hangindent=1em \textsc{Piazzi}: Lezioni elementari di astronomia \mbox{.\ .\ .\ .\ .\ .\ .\ .\ .\ .\ .\ .\ .} \hfill\mbox{1820 Jan. 6}\quad\mbox{--- 617}\par\smallskip
\noindent\hangindent=1em \textsc{Hallaschka}: Elementa eclipsium, quae potitur tellus etc. \mbox{.\ .\ .\ .\ .\ .\ .} \hfill\mbox{1820 Jan. 29}\quad\mbox{--- 619}\par\smallskip
\noindent\hangindent=1em \textsc{Bode}: Astronomisches Jahrbuch für das Jahr 1822 \mbox{.\ .\ .\ .\ .\ .\ .\ .\ .\ .} \hfill\mbox{1820 Febr. 7}\quad\mbox{--- 619}\par\smallskip
\noindent\hangindent=1em \textsc{Gerling}: Sonnenfinsterniss 1820 Sept. 7 \mbox{.\ .\ .\ .\ .\ .\ .\ .\ .\ .\ .\ .\ .} \hfill\mbox{1820 Juni 26}\quad\mbox{--- 622}\par\smallskip
\noindent\hangindent=1em \textsc{Bode}: Astronomisches Jahrbuch für das Jahr 1823 \mbox{.\ .\ .\ .\ .\ .\ .\ .\ .} \hfill\mbox{1821 Jan. 8}\quad\mbox{--- 625}\par\smallskip
\noindent\hangindent=1em \textsc{Brambilla}: Effemeridi astronomiche di Milano per l'anno 1823 \mbox{.\ .\ .\ .\ .} \hfill\mbox{1823 April 21}\quad\mbox{--- 629}\par\smallskip
\noindent\hangindent=1em \textsc{Bode}: Astronomisches Jahrbuch für das Jahr 1826 \mbox{.\ \raisebox{.25ex}{-}\ .\ .\ .\ .\ .\ .\ .\ .\ .} \hfill\mbox{1824 Jan. 3}\quad\mbox{--- 630}\par\smallskip
\noindent\hangindent=1em \textsc{Pasquich}'s Ehrenrettung \mbox{.\ .\ .\ .\ .\ .\ .\ .\ .\ .\ .\ .\hspace{.6em}\ .\ .\ .\ .\ .\ .} \hfill\mbox{1824 März 3}\quad\mbox{--- 634}\par\smallskip
\noindent\hangindent=1em Connaissance des tems etc.\ .\ .\ . pour l'an 1826 \mbox{.\ .\ .\ .\ .\ .\ .\ .\ .\ .\ .} \hfill\mbox{1824 Mai 17}\quad\mbox{--- 638}\par\smallskip
\noindent\hangindent=1em \textsc{Damoiseau}: Tables de la lune etc.\ .\ .\ . Paris 1824 \mbox{.\ .\ .\ .\ .\ .\ .\ .\ .} \hfill\mbox{1825 Febr. 7}\quad\mbox{--- 640}\par\smallskip
\noindent\hangindent=1em Observations astronomiques. faites à l'observatoire royal de Paris. Paris 1825 \hfill\mbox{1826 Aug. 14}\quad\mbox{--- 641}\par\smallskip
\noindent\hangindent=1em J. F. W. \textsc{Herschel} and \textsc{James}. \textsc{South}: Observations of 380 double and triple stars, made in the. years 1821, 1822 and 1823.\ \raisebox{.15ex}{.} London \mbox{.\ .\ .\ .\ .\ .} \hfill\mbox{1826 Dec. 18}\quad\mbox{--- 642}\par\smallskip
\noindent\hangindent=1em Astronomical observations .made at Göttingen from 1756 to 1761 by \textsc{Tobias Mayer}. Published by order of the commissioners of latitude. London 1826 \hfill\mbox{1827 Oct. 6}\quad\mbox{--- 644}\par\smallskip
\noindent\hangindent=1em Connaissance des tems etc.\ .\ .\ . pour l'an 1829 \mbox{.\ .\ .\ .\ .\ .\ .\ .\ .\ .} \hfill\mbox{1828 Jan. 10}\quad\mbox{--- 645}\par\smallskip
\noindent\hangindent=1em \textsc{Moll}: Mercurius in Sole visus. 1832 Mai 5. Utrecht. Amsterdam 1833 \mbox{.\ .} \hfill\mbox{1833 Oct. 21}\quad\mbox{--- 648}\par\smallskip
\noindent\hangindent=1em \textsc{Steinheil}: Preisschrift über zwei neue Photometer \mbox{.\ .\ .\ .\ .\ .\ .\ .\ .} \hfill\mbox{1835 März 5}\quad\mbox{--- 649}\par\smallskip
\noindent\hangindent=1em \textsc{Beer} und \textsc{Mädler}: Beiträge zur physischen Kenntniss der himmlischen Körper im Sonnensysteme. Weimar 1841. \mbox{.\ .\ .\ .\ .\ .\ .\ .\ .\ .\ .\ .}\par\smallskip
\noindent\hangindent=1em Fragments sur les corps célestes du système solaire. Paris 1840 \mbox{.\ .\ .\ .} \hfill\mbox{1841 Aug. 30}\quad\mbox{--- 652}\par\smallskip
\noindent\hangindent=1em \textsc{Beima}: De annulo Saturni. Leyden 1842 \mbox{.\ .\ .\ .\ .\ .\ .\ .\ .\ .\ .\ .\ .} \hfill\mbox{1842 Dec. 5}\quad\mbox{--- 654}\par\smallskip
\noindent\hangindent=1em \emph{Bemerkungen von Schering} \mbox{.\ .\raisebox{-.3ex}{.}\ .\ .\ .\ .\raisebox{-.3ex}{.}\ .\raisebox{-.3ex}{.}\ .\ .\ .\ \raisebox{.1ex}{\rotatebox{35}{-}}\ .\ .\ .\ .\ .\ .\raisebox{-.3ex}{.}\ .\ .\ .\ .\ .\ .\ .\ .} \hfill\mbox{}\quad\mbox{--- 654}\par\smallskip
\endgroup
\bigskip
\begin{center}
GÖTTINGEN,\par\medskip
{\small GEDRUCKT IN DER DIETERICHSCHEN UNIVERSITÄTS-DRUCKEREI}\par\smallskip
{\footnotesize W. FR. KAESTNER.}
\end{center}
